\documentclass[11pt]{article}

\usepackage[a4paper, margin=2.4cm]{geometry}
\usepackage[T1]{fontenc}
\usepackage[utf8]{inputenc}
\usepackage{lmodern}
\usepackage{amsmath, amssymb, mathtools}
\usepackage{booktabs}
\usepackage{xcolor}
\usepackage{mdframed}

\definecolor{oldcolor}{RGB}{200,50,50}
\definecolor{newcolor}{RGB}{30,120,60}

\newcommand{\stepbox}[2]{%
  \noindent\textbf{Step #1.}\quad #2\par\medskip
}

\title{LPV-SS vs.\ LPV-LFR Forward Pass\\
\large Stepwise Comparison for the Gantry Baseline}
\author{}
\date{}

\begin{document}
\maketitle

% -----------------------------------------------------------------------
\section*{Shared Setup}
% -----------------------------------------------------------------------

Both approaches start from the same mechanical model and notation.

\medskip
\noindent\textbf{State and input:}
\[
  x = \begin{bmatrix} q \\ \dot{q} \end{bmatrix} \in \mathbb{R}^6,
  \qquad
  u \in \mathbb{R}^3,
  \qquad
  Y = x_3 \quad (\text{payload Y-position, quasi-LPV scheduling variable}).
\]

\noindent\textbf{Net force:}
\[
  f_{\mathrm{net}} := -Kq - C\dot{q} + u \in \mathbb{R}^3.
\]

\noindent\textbf{Parameter-dependent inertia matrix:}
\[
  M(Y) = M_0 + Y M_1 + Y^2 M_2,
\]
where $M_0, M_1, M_2 \in \mathbb{R}^{3\times3}$ are constant matrices
derived from the physical parameters.

% -----------------------------------------------------------------------
\section{Previous Approach: Collapsed LPV-SS}
% -----------------------------------------------------------------------

The parameter-dependent inertia $M(Y)$ is inverted directly at each
time step. The model is implicitly a parameter-varying state-space
system $\dot{x} = A(Y)x + B(Y)u$; the LFR structure is never made explicit.

\medskip

\stepbox{1}{Compute $f_{\mathrm{net}} = -Kq - C\dot{q} + u$.}

\stepbox{2}{Solve the $3\times3$ linear system for the acceleration:
\[
  {\color{oldcolor} M(Y)\,a = f_{\mathrm{net}},}
  \qquad a \in \mathbb{R}^3.
\]
This requires forming $M(Y)$ at the current $Y$ and calling a
numerical linear solver at every time step.}

\stepbox{3}{Assemble $\dot{x}$ directly from $a$:
\[
  {\color{oldcolor} \dot{x} = \begin{bmatrix} \dot{q} \\ a \end{bmatrix}.}
\]
The LFR latent signals $z$ and $w$ do not exist; $M(Y)^{-1}$
is implicit.}

\medskip

\begin{mdframed}[backgroundcolor=oldcolor!8, linecolor=oldcolor, linewidth=1pt]
\textbf{Limitation.} The G matrix is never formed. The scheduling
structure $\Delta(Y)$ is never made explicit. No latent signals $z$, $w$
are available for augmentation. The computation cannot be cast as an
LFR interconnection evaluation.
\end{mdframed}

% -----------------------------------------------------------------------
\section{Current Approach: LPV-LFR with Rational Loop Solve}
% -----------------------------------------------------------------------

The model is defined as the LFR interconnection $(G,\Delta(Y))$ with
constant $G$ and explicit scheduling block $\Delta(Y) = YI_6$.
The algebraic loop $L(Y)z = C_z x + D_{zu}u$ is solved analytically
using the rational (polynomial) form of $M(Y)^{-1}$.

% -----------------------------------------------------------------------
\subsection{Rational Expression for \boldmath{$M(Y)^{-1}$}}
% -----------------------------------------------------------------------

By Cramer's rule, $M(Y)^{-1} = \mathrm{adj}(M(Y))\,/\,\det(M(Y))$.
Both numerator and denominator are polynomial in $Y$.

\medskip
\noindent\textbf{Scalar shorthands (precomputed once):}
\[
  \alpha = m_1+m_2+m_b+m_h,
  \qquad
  \beta  = \tfrac{(m_1-m_2)L_b}{2},
  \qquad
  \gamma = J_b+J_h+\tfrac{(m_1+m_2)L_b^2}{4}.
\]

\noindent\textbf{Denominator} (scalar polynomial in $Y$, always $>0$):
\begin{equation}
  d(Y) = m_h\!\left(
    \alpha\gamma - \beta^2
    + 2\beta m_h Y
    + m_h(\alpha - m_h)Y^2
  \right).
  \label{eq:dY}
\end{equation}

\noindent\textbf{Adjugate coefficient matrices} (precomputed once):
\begin{equation}
  N_0 =
  \begin{bmatrix}
    m_h\gamma      & -\beta m_h     & -\beta d\,m_h \\
    -\beta m_h     & \alpha m_h     & \alpha d\,m_h \\
    -\beta d\,m_h  & \alpha d\,m_h  & \alpha(\gamma + m_hd^2) - \beta^2
  \end{bmatrix},
  \label{eq:N0}
\end{equation}
\begin{equation}
  N_1 =
  \begin{bmatrix}
    0        & m_h^2    & d\,m_h^2 \\
    m_h^2    & 0        & 0        \\
    d\,m_h^2 & 0        & 2\beta m_h
  \end{bmatrix},
  \qquad
  N_2 =
  \begin{bmatrix}
    m_h^2 & 0 & 0 \\
    0     & 0 & 0 \\
    0     & 0 & \alpha m_h - m_h^2
  \end{bmatrix}.
  \label{eq:N1N2}
\end{equation}

\noindent\textbf{Rational inverse:}
\begin{equation}
  M(Y)^{-1} = \frac{N(Y)}{d(Y)},
  \qquad
  N(Y) = N_0 + Y N_1 + Y^2 N_2.
  \label{eq:MYinv}
\end{equation}

% -----------------------------------------------------------------------
\subsection{Forward Pass (LFR-First Ordering)}
% -----------------------------------------------------------------------

\stepbox{1}{Compute $f_{\mathrm{net}} = -Kq - C\dot{q} + u$.
\quad (Same as before.)}

\stepbox{2}{Evaluate the polynomial loop solve: acceleration $a = M(Y)^{-1}f_{\mathrm{net}}$
using the rational expression \eqref{eq:MYinv}:
\[
  {\color{newcolor}
  d(Y) = m_h\!\left(
    \alpha\gamma - \beta^2
    + 2\beta m_h Y + m_h(\alpha-m_h)Y^2
  \right),
  \qquad
  N(Y) = N_0 + YN_1 + Y^2N_2,
  }
\]
\[
  {\color{newcolor} a = \frac{N(Y)\,f_{\mathrm{net}}}{d(Y)} \in \mathbb{R}^3.}
\]
No matrix solve is required at runtime — only polynomial evaluation and
a scalar division.}

\stepbox{3}{Form the LFR latent signals explicitly:
\[
  {\color{newcolor}
  z = \begin{bmatrix} a \\ Ya \end{bmatrix} \in \mathbb{R}^6,
  \qquad
  w = \Delta(Y)\,z = Y z = \begin{bmatrix} Ya \\ Y^2 a \end{bmatrix} \in \mathbb{R}^6.
  }
\]
This is the explicit evaluation of the scheduling block $\Delta(Y) = YI_6$.}

\stepbox{4}{Compute $\dot{x}$ \emph{through} the constant G matrix — not directly from $a$:
\[
  {\color{newcolor} \dot{x} = A_x x + B_w w + B_u u,}
\]
where $A_x, B_w, B_u$ are the constant submatrices of G precomputed from $M_0^{-1}$.
The signal $w$ is causally upstream of $\dot{x}$ through $B_w$.}

\stepbox{5}{Read off output:
$y = C_y x$.}

\medskip

\begin{mdframed}[backgroundcolor=newcolor!8, linecolor=newcolor, linewidth=1pt]
\textbf{Key properties.} $G$ is an explicit constant object. $\Delta(Y)$ is
explicit. $z$ and $w$ are explicit tensors available for augmentation wiring.
The loop is solved analytically (Cramer's rule pre-applied), not numerically
at each step. The model remains in LFR form throughout — it is \emph{not}
collapsed to $A(Y)x + B(Y)u$.
\end{mdframed}

% -----------------------------------------------------------------------
\section*{Side-by-Side Comparison}
% -----------------------------------------------------------------------

\begin{center}
\begin{tabular}{lll}
\toprule
 & \textbf{LPV-SS (previous)} & \textbf{LPV-LFR (current)} \\
\midrule
Solve method      & Numerical: $M(Y)a = f_{\mathrm{net}}$         & Polynomial: $a = N(Y)f_{\mathrm{net}}/d(Y)$ \\
Solve size        & $3\times3$ system per step                     & Polynomial eval + scalar division \\
G matrix          & Never formed                                   & Constant, precomputed \\
$\Delta(Y)$       & Implicit                                       & Explicit: $YI_6$ \\
$z$, $w$          & Do not exist                                   & Explicit tensors \\
$\dot{x}$ formula & $\dot{x} = [\dot{q};\, a]$                    & $\dot{x} = A_x x + B_w w + B_u u$ \\
Numerical result  & \multicolumn{2}{c}{Identical (same physics)} \\
Augmentation      & Not possible                                   & Wire to $z$ or $w$ \\
\bottomrule
\end{tabular}
\end{center}

\end{document}
